\chapter{修辞手法}


八种常用修辞：
\begin{enumerate}
    \item \textbf{基本特征}
    \item \textbf{表达效果}
\end{enumerate}


\begin{introduction}
  \item 比喻
  \item 拟人
  \item 设问
  \item 反问
  \item 排比
  \item 反复
  \item 对偶
  \item 夸张
\end{introduction}


\section*{区别易混淆的修辞}

\subsection*{（一）比喻与拟人}

\textbf{下列句子是比喻还是拟人？}
\begin{enumerate}
    \item 每条岭都像母亲那么温柔。 \hfill \textbf{比喻}
    \item 春风温柔地抚摸着你。 \hfill \textbf{拟人}
    \item “吹面不寒杨柳风”，不错的，像母亲的手抚摸着你。 \hfill \textbf{比喻}
\end{enumerate}

\subsubsection*{区别比喻与拟人}


\begin{definition}[比喻] \label{def:int} 
明喻、暗喻、借喻
（1）对事物进行比喻：用比喻来对事物某某特征进行描绘和渲染，可使事物生动形象，具体可感，引发读者联想和想象，给人以鲜明深刻的印象，并使语言文采斐然，富有很强的感染力。
（2）对道理进行比喻：用浅显易见的事物对深奥的道理加以描述，化抽象为具体，化繁为简，帮助人们深入的理解。并使语言生动形象，富有文采。
\end{definition}


\begin{definition}[拟人] \label{def:int} 
用拟人化手法，使抽象事物具体化，使具体事物人格化，生动形象地表达了某某感情，富有强烈的感染力。
\end{definition}

\begin{center}
    \renewcommand{\arraystretch}{1.8}
    \begin{tabular}{|>{\centering\arraybackslash}p{3cm}|>{\centering\arraybackslash}p{5.5cm}|>{\centering\arraybackslash}p{5.5cm}|}
        \hline
        & \textbf{比喻} & \textbf{拟人} \\
        \hline
        \textbf{核心特征} & 两物有 \textcolor{red}{相似点} & 物即人（将物当人来写） \\
        \hline
        \textbf{句中标志} & 句中有 \textcolor{red}{喻体} & 句中 \textcolor{red}{无“人”}（不出现人称词） \\
        \hline
        \textbf{表达效果} & 生动形象地表现事物的特点 & 把物人格化，更具体生动 \\
        \hline
    \end{tabular}
\end{center}

\subsubsection*{练一练}

\begin{enumerate}
    \item “像柳絮一般的雪，像芦花一般的雪，像蒲公英的带绒毛的种子一般的雪，在风中飞舞。” \\
    \textbf{理解：} 运用了拟人、排比的修辞方法，生动形象地描写了雪花洁白的颜色和漫天飞舞的美态。 \hfill ×

    \item 天边几颗调皮的星星，时隐时现，似乎在不知疲倦地和月亮捉着迷藏。 \\
    \textbf{理解：} 这句话运用拟人的修辞方法，生动地写出了星星时隐时现的情态。 \hfill √
\end{enumerate}

\subsection*{（二）区别反问与设问}

\textbf{下列句子是反问还是设问？}
\begin{enumerate}
    \item 做工苦，难道不做工就不苦吗？ \hfill \textbf{反问}
    \item 不愿意做逃得了吗？到底不能。 \hfill \textbf{设问}
\end{enumerate}

\subsubsection*{区别反问与设问}

\begin{center}
    \renewcommand{\arraystretch}{1.8}
    \begin{tabular}{|>{\centering\arraybackslash}p{3cm}|>{\centering\arraybackslash}p{5.5cm}|>{\centering\arraybackslash}p{5.5cm}|}
        \hline
        & \textbf{反问} & \textbf{设问} \\
        \hline
        \textbf{核心特征} & \textcolor{red}{答在问中}（答案隐含在问句里） & \textcolor{red}{自问自答}（先问后答） \\
        \hline
        \textbf{表达效果} & 加强语气，感情强烈 & 提醒注意，引起思考 \\
        \hline
    \end{tabular}
\end{center}

\vspace{0.5cm}
\begin{center}
    \textcolor{red}{注：一般疑问句“有问无答”，不是修辞。}
\end{center}

\subsubsection*{判断下列句子是设问还是反问？}

\begin{enumerate}
    \item 友谊是什么？友谊是联系情感的纽带，是相互沟通的桥梁，是化解矛盾的良药。 \hfill \textbf{设问}
    \item 你不觉得我们的战士是可爱的吗？你不以我们的祖国有这样的英雄而自豪吗？ \hfill \textbf{反问}
    \item 这是多么深刻的教诲啊！难道不值得我们永远牢记吗？ \hfill \textbf{反问}
    \item 你喜欢上语文课吗？喜欢。 \hfill \textbf{一般疑问句}
\end{enumerate}

\subsubsection*{练一练}

\begin{enumerate}
    \item 花儿为什么这样红？首先有它的物质基础。 \\
    \textbf{理解：} 运用了设问的修辞方法，故意先提出问题，然后自己回答，提醒人们思考。突出了物质基础是花儿红的首要原因。 \hfill √

    \item 天气这么冷，你怎么能穿这么少的衣服呢？快把我这件大衣披上吧。 \\
    \textbf{理解：} 这句话运用设问的修辞方法，强调了说话人对对方的关心、照顾、体贴。 \hfill ×
\end{enumerate}

\subsection*{（三）区别排比与反复}

\textbf{下列句子是排比还是反复？}


\begin{enumerate}
    \item 沉默呵，沉默！不在沉默中爆发，就在沉默中灭亡。 \hfill \textbf{反复}
    \item 延安的歌声，是革命的歌声，战斗的歌声，劳动的歌声，极为广泛的群众的歌声。 \hfill \textbf{排比}
\end{enumerate}

\subsubsection*{区别排比和反复}

\begin{definition}[排比] \label{def:int} 
\begin{itemize}
    \item \textbf{排比说理：} 使文章条理清晰，论述详尽，透彻严密，语气强烈，无可辩驳。
    \item \textbf{排比抒情：} 使文章将感情抒发得淋漓尽致，并使文章音节铿锵，语势得到增强。
    \item \textbf{叙事排比：} 使文章深刻细致，层次清楚，表述全面，有一气呵成之感。
\end{itemize}
\end{definition}


\begin{definition}[反复] \label{def:int} 
\begin{itemize}
    \item \textbf{连续反复} 
    \item \textbf{间隔反复} 
\end{itemize}
\end{definition}

\begin{center}
    \renewcommand{\arraystretch}{1.8}
    \begin{tabular}{|>{\centering\arraybackslash}p{3cm}|>{\centering\arraybackslash}p{5.5cm}|>{\centering\arraybackslash}p{5.5cm}|}
        \hline
        & \textbf{排比} & \textbf{反复} \\
        \hline
        \textbf{结构数量} & \textcolor{red}{三个或三个以上}（结构相似或相同的分句） & \textcolor{red}{一个}（同一个词语或句子） \\
        \hline
        \textbf{重复方式} & 相似或相同的结构并列 & 被两次或多次重复 \\
        \hline
        \textbf{表达效果} & 加强语势 & 强调某种意思或感情 \\
        \hline
    \end{tabular}
\end{center}

\subsubsection*{练一练}

\begin{enumerate}
    \item 风！你咆哮吧！咆哮吧！尽力地咆哮吧！ \\
    \textbf{理解：} 这句话运用了“反复”，使用了三个“咆哮吧”，突出强调了屈原当时满腔的愤怒和对黑暗势力的憎恨。 \hfill √

    \item 那沉甸甸的稻谷，像一垄垄金黄的珍珠； \\
    炸蕾吐絮的棉花，像一箱箱雪白的珍珠； \\
    婆娑起舞的莲蓬，却又像一盘盘碧绿的珍珠。 \\
    \textbf{理解：} 这句话运用了“反复、排比”，“像……珍珠”重复使用，突出了表现了稻谷、棉花、莲蓬的特点；三句并列排比，增强语势表达了喜爱和赞叹之情。 \hfill ×
\end{enumerate}


\subsection*{夸张}

用言过其实的方法，突出事物的本质，或加强作者的某种感情，烘托气氛，引起读者的联想。


\section*{对偶}

使文章节奏刚强，看起来醒目，读起来顺口，听起来悦耳，便于传颂记忆。